\section{Recovering both initial velocities from orbit snapshots}
\label{sec:both-velocities}
The original record supplies one component of the
initial velocity. When both are withheld, the two
null-ray transports act on a vector rather than on
one initial direction. A closed sequence of ray shifts
then gives a matrix condition for stable recovery.
This section concerns two full orbit snapshots and one
initial metric anchor. It does not give a negative
result when a full metric potential is supplied as an
additional spatial field.

Let the missing velocity variation be
$f\in H^s(\mathbb T;\mathbb R^2)$ in the chosen
orthonormal frame. The initial orbit position is fixed.
Replace the columns $A_\pm$ of
\eqref{eq:transport-leading} by the full parallel
transport matrices on the two null rays, and let $h$
be the componentwise mean-zero primitive of $f$.
The same leading formula and compact remainder hold.
After rotating the terminal output by $A_-^{-1}$,
the principal operator on $H(\theta)=h(\theta+d)$ is
\[
 I-\mathcal U_d,\qquad
 (\mathcal U_d H)(\theta)
  =W_d(\theta)H(\theta-2d),\qquad
 W_d=A_-^{-1}A_+\in \mathrm{SO}(2).
\]
This is a unitary weighted rotation in $L^2$.
An equivalent $H^{s+1}$ conjugation changes the
operator by a compact term because the coefficient
matrix is smooth.

\begin{proposition}[A rational ray-holonomy test]
\label{prop:both-ray-holonomy}
For irrational $d/\pi$ the linearized two-orbit-snapshot
map cannot have an energy-norm lower bound for both
missing velocity components, even after excluding
finitely many directions.
For rational $d/\pi=p/q$ in lowest terms define
\[
 \mathcal H_d(\theta)=
     \prod_{j=0}^{q-1} W_d(\theta-2jd) .
\]
The principal operator is bounded below precisely when
$I-\mathcal H_d$ is uniformly invertible.
In that case the exact observation derivative has a
lower bound up to an $H^{s-1}$ velocity remainder.
A sufficiently early finite set of extra low-mode
orbit observations gives a bounded-below derivative
and a local nonlinear Lipschitz inverse.
\end{proposition}

\begin{proof}
For irrational rotation, conjugating $\mathcal U_d$
by multiplication by $e^{im\theta}$ on the
complexified space multiplies it by $e^{-2imd}$.
Its nonempty essential spectrum is therefore the
whole unit circle. Since the operator is unitary,
the point $1$ has a weakly null approximate sequence.
Conjugation to the Sobolev norm and the smoothing
observation remainder do not remove this sequence.
The mean and the linearized integral momentum
constraint amount to finitely many restrictions,
also harmless for the weakly null sequence.

For rational rotation $\mathcal U_d^q$ is
multiplication by $\mathcal H_d$.
If $I-\mathcal H_d$ is uniformly invertible,
factor $I-\mathcal U_d^q$ to obtain the claimed
principal lower bound:
\[
 I-\mathcal U_d^q
    =(I+\mathcal U_d+\cdots+\mathcal U_d^{q-1})(I-\mathcal U_d).
\]
If $\mathcal H_d(\theta_0)$ has eigenvalue one, use
an envelope supported near $\theta_0$ and a vector
fixed by that holonomy. Sum its first $q$ weighted
translates; they have disjoint supports for small
enough envelopes, and the residual telescopes to
$(I-\mathcal U_d^q)$ applied to the original envelope.
It becomes arbitrarily small as the support shrinks.
Attach a high Fourier carrier with frequency a
multiple of $q$ to commute with each shift.
Taking account of the Sobolev normalization and
the compact remainders by choosing arbitrarily
large carriers shows failure of an energy lower
bound even after compact corrections.
The vector energy estimate and the earlier low-mode
argument of Theorem~\ref{th:fixed-rank-two} prove
the final claim. The local nonlinear argument
is unchanged when both initial velocities vary
subject to the affine momentum constraint.
\end{proof}

\begin{theorem}
\label{th:both-circular}
For a fixed sufficiently short rational delay,
there are circular-loop Gowdy backgrounds, and open
neighborhoods containing general nearby Cauchy data,
on which two full orbit snapshots are locally
bi-Lipschitz for both missing initial velocities.
The construction works for all sufficiently small
nonzero initial radii outside a countable set.
\end{theorem}

\begin{proof}
Take a circular-loop solution with angular winding
$k_0\ge1$, initial radius $\rho>0$, and zero initial
radial velocity. The chosen interval is short
enough that the radius stays positive for small
$\rho$.
In the radial and unit angular frame,
the spatial connection is $k_0\cosh R(t)$ times
the generator of rotations, and the temporal
connection vanishes. Hence the angle of the
holonomy in Proposition~\ref{prop:both-ray-holonomy}
is
\[
 -2k_0 q\int_a^c\cosh R(t)\,dt
   \equiv -2k_0q\int_a^c(\cosh R(t)-1)\,dt
                   \pmod{2\pi}.
\]
For any sufficiently small $\rho>0$ the
last integral is positive and the angle
is strictly between $-2\pi$ and zero.
The principal lower bound therefore holds.

In this frame all coefficients of the exact
linearized equation are independent of $\theta$,
so there is a $2\times2$ terminal transfer matrix
on each complex Fourier mode.
These matrices depend analytically on $\rho$,
including $\rho=0$ by using the rotating
orthonormal frame at the target center.
At $\rho=0$ the equation is the constant-background
scalar Gowdy equation on each physical component.
A Fourier mode $k$ written in the rotating frame
has the two shifts $k\pm k_0$ in the fixed frame.
Choose the rational delay as in
Proposition~\ref{prop:safe-gap}, so every one of
those scalar transfers, including zero frequency,
is nonzero. For each fixed $k$ the transfer
determinant is analytic and nonzero at
$\rho=0$. The union of its zero sets over all
integer $k$ is countable.
Away from this set the derivative has no kernel.
The lower bound up to compact remainder then
gives a full energy lower bound, by the same
compactness argument as before.
Continuity of the derivative in the Cauchy
data proves stability in nearby fully nonlinear
solutions.
\end{proof}


\paragraph{The holonomy of the Ernst pair.}
At a short rational delay $d=\pi/L$ with $L>2k_0$,
the rank-two Ernst pair also satisfies the principal
test for both velocities when its amplitude is
small and nonzero.
Its relative ray rotation is $-2B(w(c,\theta))$.
Use $B(w)=q w-(qg/2)w^2+O(w^3)$.
Summing the first harmonic over the $L$ ray shifts
gives zero, while the sum of its square is $L/2$.
The resulting holonomy angle is therefore
\[
 \frac{qgL}{2}\,\epsilon^2u_{k_0}(c)^2+O(\epsilon^3),
\]
uniformly in $\theta$ and with the sign of the
chosen geodesic branch. It is nonzero for small
nonzero amplitude and $\;gq\ne0$.
One full terminal orbit snapshot consequently
gives the high-frequency estimate in this family.
A finite set of early low-mode coefficients
removes any finite-dimensional kernel.
The next theorem gives explicit two-terminal
placements that do not require locating this kernel.


\paragraph{The local two-time compatibility equation.}
At the circular-loop backgrounds of
Theorem~\ref{th:both-circular}, the bounded-below
two-component derivative is also onto the two
terminal orbit fields. To see this, extend the
mean-zero primitive to an isomorphism from $H^s$
to $H^{s+1}$ by adding the mean mode.
Its difference from the leading expression is
finite rank. The rational holonomy condition
makes its composition with $I-\mathcal U_d$
invertible. The exact observation derivative is
a compact perturbation of this invertible
operator, hence has index zero.
Its trivial kernel then gives surjectivity.

It follows from the inverse function theorem that
the reduced wave-map problem with prescribed orbit
positions on the two slices has a locally smooth
solution for initial positions and terminal values
in their natural Sobolev neighborhoods.
The periodic vacuum metric exists with a specified
anchor exactly on the image of the momentum condition
\[
 \int_{\mathbb T}\langle U_t(a),U_\theta(a)\rangle\,d\theta=0 .
\]
When $U_\theta(a)$ is nonzero somewhere, this is
locally one compatibility equation on the measured
orbit positions. Indeed its derivative in the recovered
velocity is a nonzero linear functional, and the
two-snapshot reconstruction is a local diffeomorphism
for the reduced equations. At the circular-loop
reference with the initial orbit position held fixed,
the linearized compatibility condition is that the
mean angular component of the terminal displacement
in the rotating frame vanish. The initial angular
velocity enters the constraint only through its
zero Fourier mode, and that mode is decoupled
with nonzero scalar transfer in this frame.

\subsection{Three orbit snapshots in the exact rank-two family}

\begin{theorem}
\label{th:both-ernst}
At each sufficiently small nonzero amplitude of the
single-mode rank-two Gowdy-to-Ernst pair, two suitably
placed later full orbit snapshots together with the
initial full orbit snapshot recover \emph{both} missing
initial velocities locally with a Lipschitz estimate
in the wave-energy Sobolev norms. The placements
can be chosen as different short rational gaps,
independently of the amplitude.
\end{theorem}

\begin{proof}
Let the two delays be $d_i=\pi/L_i$, $i=1,2$,
with both integers $L_i$ sufficiently large,
distinct, and greater than $2k_0$.
At the zero-amplitude vertical background, the two
time-scaled linearized channels have inverse-square
potentials $1/(4t^2)$ and $-3/(4t^2)$.
At these short rational placements neither scalar
channel has a finite-frequency kernel.
For the negative potential the Pr\"ufer correction
to $kd_i$ has the opposite sign from the
positive-potential case and absolute magnitude
less than $\pi/L_i$ for large $k$.
For remaining modes choose $L_i$ large so that
$\varphi(c_i)<\pi$; the phase remains strictly
positive since its derivative is $k$ at each
zero of $\varphi$. The zero spatial mode has
nonzero Jacobi transfers.
On high frequencies both channels have leading
normalized gain $\sin(kd_i)$.

Write $\mathcal A_\epsilon$ for the joint derivative
with respect to the unknown two-component velocity
and $\epsilon$ for the seed amplitude.
It has the norm expansion
$\mathcal A_\epsilon=\mathcal A_0+
\epsilon \mathcal A_1+O(\epsilon^2)$.
Let $L$ be the least common multiple of $L_1,L_2$.
Split the velocity into $f_B$, its high modes
divisible by $L$, and $f_G$, the remaining modes.
The free joint map $\mathcal A_0$ is bounded below on
$G$, since all finite transfers are nonzero and
each large noncommon multiple has a nonzero free
gain at one of the times.

At each output Fourier mode, project the four
observations orthogonally away from the two
columns of $\mathcal A_0$ there, in the
energy-normalized coordinates. Call this bounded
projection $\Pi^\perp$.
For an input high common multiple $k$, the first
variation of the transport angles maps to sidebands
$k\pm k_0$. Up to the common rotation of the
velocity vector and Fourier normalization, its
two time coefficients at each sideband are
\[
 \big((-1)^{k/L_1}u_{k_0}(a+d_1),
      (-1)^{k/L_2}u_{k_0}(a+d_2)\big).
\]
Indeed the two shifted primitives coincide at each
common resonance and the rotations change with
opposite signs on the two rays.
The leading free column at that sideband has time
coefficients proportional to
\[
 \big((-1)^{k/L_1}\sin(k_0d_1),
      (-1)^{k/L_2}\sin(k_0d_2)\big).
\]
The two vectors are not collinear: near zero,
\[
 \frac{u_{k_0}(a+d)}{\sin(k_0d)}
    =\frac1{a k_0}-\frac{d}{2a^2 k_0}+O(d^2)
\]
is strictly decreasing. The relative signs at
the two sidebands do not affect the projected gain.
Distinct high common multiples give disjoint
sidebands because $L>2k_0$.
The exact free projections at those sidebands
approach the displayed limiting projection with
error $O(1/|k|)$, uniformly along the resonant
subsequence. The differentiated transport remainder
is compact as before. Consequently
$\Pi^\perp\mathcal A_1$ is bounded below on $B$
after increasing the finite cutoff if needed.

The good-mode estimate and the projected estimate
therefore give
\begin{align*}
 \|f_G\|_{H^s}
   &\le C\big(\|\mathcal A_\epsilon f\|+
                 |\epsilon|\|f\|_{H^s}\big),\\
 \|f_B\|_{H^s}
   &\le C'\big(\|\mathcal A_\epsilon f\|/|\epsilon|+
                  \|f_G\|_{H^s}+|\epsilon|\|f\|_{H^s}\big).
\end{align*}
For small nonzero amplitude the error terms
are absorbed. The exact derivative is then bounded
below, and continuity of the nonlinear derivative
gives the local result.
\end{proof}
